\section{引言：非形式化与形式化数学推理的桥梁}

本讲由 CMU 的 Sean Welleck 主讲，探讨如何利用 AI 弥合\textbf{非形式化数学}（informal math）与\textbf{形式化数学}（formal math）之间的鸿沟，结合两者各自的优势。

\begin{importantbox}{两种数学表示的对比}
\textbf{非形式化数学}：以自然语言（文本、LaTeX、图像等）表示，极为灵活，但无法保证正确性。LLM 在中间步骤中可能犯微妙错误。\\
\textbf{形式化数学}：以编程语言（Lean、Coq、Isabelle）编写，代码可编译 = 证明正确。提供绝对的正确性保证，但编写门槛高、过程繁琐。
\end{importantbox}

\subsection{形式化证明助手（Lean）简介}

Lean 是一种交互式证明助手（Interactive Theorem Prover），核心概念：
\begin{itemize}
    \item \textbf{定理声明}：用形式化语言表述需要证明的命题
    \item \textbf{证明状态}（Proof State）：当前上下文（假设）和待证目标
    \item \textbf{策略}（Tactic）：每一步应用的推理规则，改变证明状态
    \item \textbf{QED}：所有目标消解后证明完成
\end{itemize}

\begin{knowledgebox}{Mathlib：形式化数学的基础设施}
Mathlib 是 Lean 的数学库项目，众多贡献者共同构建数学基础定义和定理。它使形式化证明可以复用已有结果，数学家 Terence Tao 已成功将其论文定理在 Lean 中形式化。
\end{knowledgebox}

\subsection{形式化数学对 AI 的重要性}

\begin{enumerate}
    \item \textbf{输出可验证}：防止不正确的数学推理
    \item \textbf{完美奖励信号}：可用于 RL 训练——证明编译通过即正确
    \item \textbf{推理测试平台}：从简单到任意困难的推理问题
\end{enumerate}

形式化定理证明近年进展迅速：从 GPT-f（2020）到 DeepSeek-Prover，mini-F2F 基准测试成绩持续飙升。

